% ===================================================================== ANALYSIS
\section{Analysis}
\label{sec:analysis}

\subsection{Assumptions}
\begin{assumption}[Engine correctness]
\label{as:engines}
Each engine conforms to \cref{def:engine}. $\Csrc$ is safe under its admitted source quorum.
$\Ctgt$ is always safe and, after GST, live under $q_2$ with $2q_2-n>f$ and $q_2\le n-f$.
\end{assumption}

\begin{assumption}[Deterministic execution]
\label{as:exec}
$\delta_{\mathrm{exec}}$ is deterministic on well-formed blocks: correct validators that apply the same
sequence from the same state compute identical roots.
\end{assumption}

\begin{assumption}[Metadata-independent reversible window]
\label{as:metadata}
For transactions finalized in $[\hs,\hr)$, execution is a function only of the ordered transactions
and the parent state. It excludes engine-specific metadata and live external inputs that may differ
between the target attempt and legacy replay.
\end{assumption}

\begin{assumption}[Cryptography]
\label{as:crypto}
$\mathrm{root}(\cdot)$ and $\Hash$ are collision resistant, signatures are EUF-CMA secure, and the
adversary controls at most $f$ keys, with $n\ge3f+1$.
\end{assumption}

We write $\mathrm{committed}_i(h,B)$ when correct validator $i$ holds a finalization certificate for
$B$ at height $h$. Two correct validators \emph{conflict} at $h$ if their certified blocks have
different execution-state roots. Engine-tagged block hashes alone do not define a cross-engine conflict.
Throughout, $|S\cap T|\ge|S|+|T|-n$ for any $S,T\subseteq\VM$, and a set of more than $f$ validators
contains a correct one. Every quorum statement below follows from these two facts.

\subsection{Boundary certificate}
\begin{lemma}[Quorum intersection]
\label{lem:qi}
Any two $(n-f)$-subsets of $\VM$ intersect in at least $n-2f\ge f+1$ validators and therefore share a
correct validator.
\end{lemma}

\begin{theorem}[Safe and available threshold window]
\label{th:window}
Let a fixed-candidate boundary certificate require $q$ matching attestations. The certificate is
\emph{available} (formable after GST whenever all correct validators are ready on the same body and
$f$ validators stay silent) if and only if $q\le n-f$. It is \emph{intersection-safe} (any two
certificates share a correct signer) if and only if $2q-n>f$. Hence both hold if and only if
\begin{equation}
  q\in\Bigl(\frac{n+f}{2},\;n-f\Bigr]\cap\mathbb Z ,
  \label{eq:window}
\end{equation}
a set that is non-empty exactly when $n\ge3f+1$. In particular, $2f+1$ lies in \eqref{eq:window}
if and only if $n=3f+1$, and $q=n-f$ maximizes the guaranteed correct intersection $2q-n-f=n-3f$.
\end{theorem}
\begin{proof}
\emph{Availability.} If $q>n-f$, the $f$ silent validators leave at most $n-f<q$ signers, so no
certificate forms. If $q\le n-f$, the $n-f$ correct ready validators deliver $q$ matching shares after
GST.
\emph{Intersection.} Two $q$-subsets can be chosen with intersection exactly $\max(0,2q-n)$. If
$2q-n\le f$, place the whole intersection inside the corrupted set; the two certificates then share no
correct signer. If $2q-n>f$, every intersection has more than $f$ members and contains a correct one.
\emph{Window.} Combining the two gives \eqref{eq:window}. The integer $n-f$ lies in the window if and
only if $2(n-f)-n>f$, that is $n>3f$. For $q=2f+1$, the lower bound requires $4f+2-n>f$, i.e.,
$n<3f+2$, and with $n\ge3f+1$ this forces $n=3f+1$.
\end{proof}

\begin{theorem}[Decision uniqueness]
\label{th:unique}
Under the one-candidate-body rule of \cref{def:cutcert}, at most one tuple $\beta$ is certified by a
valid \CutCert{} in epoch $e$. Every correct validator that accepts a \CutCert{} therefore accepts the
same $\hs$, parent hash, and boundary root.
\end{theorem}
\begin{proof}
Let $\pi$ and $\pi'$ certify $\beta$ and $\beta'$. By \cref{lem:qi} their signer sets share a correct
validator. That validator signs only its locked tuple in epoch $e$, so $\beta=\beta'$ unless a
signature is forged, which \cref{as:crypto} excludes.
\end{proof}

\subsection{Authority transfer}
The following theorem isolates a fence-only contract: correct fence signers have not authorized
post-boundary source progress and install their prohibition before sharing. It does not assume the
readiness lock. Its converse concerns executions admitted by that contract without a different
source-retirement mechanism, not the locked \sage{} design.
\begin{theorem}[Cross-engine authority exclusion under partial delivery]
\label{th:authority-exclusion}
Let $F$ be the signer set of a valid \FenceCert{} with $|F|=q_F$, and let $L$ be a quorum of size
$q_1$ under the source engine's native finality policy. Suppose at most $f$ validators are Byzantine,
$q_F+q_1>n+f$, and every correct fence signer (i)~has issued no source authorization at any
$h\ge\hs$ and (ii)~durably disables such authorization before releasing its share. Then no source
quorum finalizes at any $h\ge\hs$. This holds even if the \FenceCert{} reaches only a target quorum.
Conversely, if $q_F+q_1\le n+f$, there is an execution with partial delivery in which a valid
\FenceCert{} activates the target and a source quorum still finalizes at $\hs$.
\end{theorem}
\begin{proof}
$|F\cap L|\ge q_F+q_1-n>f$, so $F\cap L$ contains a correct validator $v$. By (i), $v$ contributed no
source authorization at $h\ge\hs$ before signing, and by (ii) and \cref{def:fencecert} it can
contribute none afterwards. $L$ therefore lacks a required signature and cannot form. The argument uses
only the signers' install-before-share order and the intersection bound. It never assumes that any
other validator received the certificate.

For the converse, let $k=\max(0,q_F+q_1-n)\le f$. Choose $F$ and $L$ with $|F|=q_F$, $|L|=q_1$ and
$|F\cap L|=k$, which is possible because $q_F+q_1-k\le n$, and corrupt exactly $F\cap L$. The correct
members of $F$ fence and release shares, so the \FenceCert{} forms and the target activates. The
adversary withholds the certificate from $L\setminus F$, whose correct members therefore never fence.
The Byzantine members of $F\cap L$ also sign in the source (\cref{sec:model}), so $L$ finalizes at
$\hs$.
\end{proof}

\begin{theorem}[Fence admissibility]
\label{th:fence-admission}
For a source policy with quorum $q_1$, there exists a fence quorum $q_F$ that is both available
($q_F\le n-f$) and excluding ($q_F+q_1>n+f$) if and only if $q_1>2f$. When it exists, $q_F=n-f$ is
admissible, and the admissible set is the integer interval $(n+f-q_1,\;n-f]$, which contains $q_1-2f$
values.
\end{theorem}
\begin{proof}
An admissible $q_F$ exists if and only if $n+f-q_1<n-f$, i.e., $q_1>2f$. The interval
$(n+f-q_1,n-f]$ contains exactly $(n-f)-(n+f-q_1)=q_1-2f$ integers, and its maximum is $n-f$.
\end{proof}

\begin{corollary}[Fence-only majority-source counterexample]
\label{cor:cross-split}
Let $n=7$, $f=2$, $q_F=5$, and let the source finalize with a majority $q_1=4$. Then $q_F+q_1=n+f$,
so the premise of \cref{th:authority-exclusion} fails. A five-validator target side and a
four-validator source side can overlap in exactly the two Byzantine validators and both finalize at
$\hs$. Fencing only the \FenceCert{} signers cannot prevent this fence-only execution. The policy
$q_1=2f$ is rejected by \cref{th:fence-admission}; this is not an admitted locked execution.
\end{corollary}
\begin{proof}
This is the converse of \cref{th:authority-exclusion} with $k=2$. Label the validators
$1,\ldots,7$, with $\{4,5\}$ Byzantine. Let $F=\{1,2,3,4,5\}$ fence and activate the
target, and let $L=\{4,5,6,7\}$ finalize in the source; correct validators $6$ and $7$ have not
installed a fence. The two sets share only $\{4,5\}$, which sign on both sides by the cross-domain
capability of \cref{sec:model}. Both sides therefore commit at $\hs$ in the fence-only contract.
\end{proof}

For a majority source $q_1=\lfloor n/2\rfloor+1$, \cref{th:fence-admission} reduces to $n\ge4f$.
Hence $(n,f)=(6,1)$ is admissible, while $(7,2)$ and $(20,6)$ are not. \Cref{fig:quorum} visualizes
the boundary and fence-only conditions. The arithmetic example below does not establish that the
fence certificate is independently necessary after a durable readiness lock.

\begin{proposition}[Independence of the two inequalities]
\label{prop:independence}
Neither the boundary window \eqref{eq:window} nor the fence-only admission inequality
\eqref{eq:intro-fence} implies the other as an arithmetic condition. (a)~At
$(n,f,q,q_1,q_F)=(7,2,5,4,5)$ the \CutCert{} threshold is safe and available, but the fence-only
exclusion inequality fails. (b)~At $(n,f,q,q_1,q_F)=(6,1,3,5,5)$ that exclusion inequality holds
($10>7$), but the boundary threshold is unsafe ($2\cdot3-6=0\le1$).
\end{proposition}
\begin{proof}
Direct substitution into \eqref{eq:window} and \eqref{eq:intro-fence}, using \cref{cor:cross-split}
for (a).
\end{proof}


\begin{lemma}[Source quiescence and fence liveness]
\label{lem:fence-live}
Assume \cref{rule:lock}, a valid \CutCert{} on $\beta$, and an admitted source policy with $q_1>2f$.
(i)~No source quorum finalizes at any height $h>\hb$ of generation $g_1$. (ii)~If, in addition, every
correct validator is eligible to attest and ready on $\beta$ and, after GST, holds $\pi_{\mathrm B}$
and has durably adopted the same certified $\mathcal M_{\mathrm B}$, then all correct validators
can supply fence shares. If an available collector receives them within $\Delta$ of the last
durable fence write, it can form a \FenceCert{} with any admissible $q_F\le n-f$ from those
shares; its local verification time is additional.
\end{lemma}
\begin{proof}
(i)~$\pi_{\mathrm B}$ has $n-f$ signers, at least $n-2f$ of them correct; call this set $A$. Every
validator outside $A$ is either Byzantine or a correct non-signer, so $|\VM\setminus A|\le 2f<q_1$.
Any source quorum above $\hb$ therefore contains a member of $A$. By \cref{rule:lock}, that member
attested before any source authorization above $\hb$ and cannot authorize one afterwards. Its
prohibition is committed before its attestation becomes visible: a crash before the commit
releases no attestation, while one afterwards reloads the prohibition before source use. Thus
restart preserves the abstract invariant if the durable store meets \cref{rule:lock}; no runtime
crash result is inferred. Delayed signatures cannot form a source quorum above $\hb$: every
such quorum needs an authorization a member of $A$ never issued.
(ii)~Eligibility means that each correct validator has issued no source authorization above $\hb$.
Given common readiness, it attests under \cref{rule:lock} before any such authorization; the rule
then prohibits later ones. All $n-f$ correct validators satisfy the no-prior-authorization
condition of \cref{def:fencecert}, install the fence after manifest adoption, and release a share.
Those $n-f\ge q_F$ shares suffice at the collector under the stated delivery bound.
\end{proof}

\begin{remark}
Part~(i) is sufficient for source exclusion in the admitted locked design: no \FenceCert{} is used
in its proof. Removing only that certificate, while preserving the lock, $q=n-f$, and $q_1>2f$,
cannot create residual source finalization. The counterexample in \cref{cor:cross-split} does not
satisfy these premises. Part~(ii) requires eligibility as well as readiness: a correct validator
that voted above $\hb$ before learning of $\beta$ cannot later sign a fence share. In the worst
case only the $n-2f$ correct attesters remain eligible; correct-only formation then requires
$q_F\le n-2f$, compatible with \cref{th:fence-admission} only when $q_1>3f$. Without the
attestation prohibition, fence-only exclusion still follows from \cref{th:authority-exclusion}
once a valid \FenceCert{} forms, but such prior votes may leave the boundary stale.
\end{remark}

\begin{proposition}[Policy-bound retirement evidence]
\label{prop:retirement-evidence}
Under \cref{as:crypto,def:fencecert}, a valid \FenceCert{} with $q_F$ distinct shares attests that
at least $q_F-f$ correct signers durably adopted the same complete manifest and installed the
specified generation fence before sharing. A \CutCert{} alone does not attest these later actions.
This is an evidence property, not additional source-exclusion safety under
\cref{lem:fence-live}(i), and it does not imply adoption by every validator.
\end{proposition}
\begin{proof}
At most $f$ shares belong to Byzantine validators. Each remaining signer releases its share only
after the manifest and fence writes, and the signed identity binds the manifest, boundary,
generation, and quorum policies. Signature unforgeability and collision resistance preserve this
binding. Conversely, an execution may form a \CutCert{} and delay all manifest messages; its correct
attesters remain locked but have not adopted a manifest or installed its policy-bound fence. Thus
the boundary certificate does not certify those writes. A remote verifier still trusts correct
signers' specified storage behaviour; a signature is not hardware attestation of an actual flush.
\end{proof}

\paragraph{Alternative source-adapter contract}
Retirement need not occur at readiness in every possible adapter. The following conditional
alternative does not replace \cref{rule:lock} in the protocol above. It instead asks whether a
source-finalized control block can itself retire the source before a later activation certificate.

\begin{proposition}[Retirement under final-parent mediation]
\label{prop:final-parent-retirement}
Fix a source generation $g$ with at most $f$ Byzantine validators. Assume native source safety:
valid finality evidence authenticates a unique parent-linked finalized prefix, and every accepted
source certificate contains $q_1>f$ distinct authorized signers. The condition $q_1>f$ alone does not establish source safety.
Suppose every correct validator, before creating any externally releasable source authorization,
adopts authenticated finality for the exact finalized parent at the preceding height.
Adopting a designated Boundary control block $B$ at $\hb$ atomically and durably installs a stop on
all source authorization above $\hb$ in $g$, before signing services can use $B$ as a parent.
All proposal, vote, catch-up and buffered-release paths serialize against that record and recheck
it before release. Restart reloads and validates the stop before enabling those paths; ambiguous
recovery refuses. If $B$ finalizes, no source block above $\hb$ finalizes in $g$.
This conclusion does not require an activation certificate.
\end{proposition}
\begin{proof}
Suppose a later source block finalizes in $g$. Prefix uniqueness gives a finalized first child
$C$ of $B$ at $\hb+1$. Its certificate contains a correct signer $v$, since $q_1>f$.
Before authorizing $C$, $v$ had to adopt its exact finalized parent $B$; that adoption installed
the stop before $B$ became usable by signing services. Thus $v$ could not authorize $C$,
contradicting the certificate. An earlier prepared or buffered authorization cannot bypass this
order: its creation also requires adoption of $B$. Using a different finalized parent at $\hb$
contradicts prefix uniqueness. A validator that has not adopted $B$ cannot authorize its child,
and correct restart preserves the prohibition. The signer $v$ need not belong to the Boundary signer set:
the adapter rule applies to every correct signer, including non-recipients until they adopt $B$.
\end{proof}

\begin{remark}[Scope of the alternative]
Under these stronger premises, a later certificate can remain a target-authorization and
prerequisite-evidence gate without being independently necessary for source retirement.
The fence-only inequality of \cref{th:authority-exclusion} is not a necessity result for this stronger adapter;
native source-safety and governance admission requirements remain separate.
Readiness need not install this stop, allowing the source to finalize $B$.
A losing Boundary vote alone does not retire the source: a valid ordinary winner may permit
its next child while retaining the losing per-height vote lock. This is not a liveness result.
The final-parent restriction is not established for the evaluated implementation;
neither the reported timings nor the existing bounded models validate this alternative.
\end{remark}

\subsection{Handoff safety and completion}
\begin{lemma}[Single finalizer during dual-run]
\label{lem:single}
For every $h\in[\hd,\hs)$, only $\Csrc$ finalizes, and it finalizes at most one block at $h$.
\end{lemma}
\begin{proof}
In \phase{dual}, $\Auth=\Csrc^{g_1}$ and the shadow only validates finalized blocks
(\cref{def:shadow}). \Cref{as:engines} gives uniqueness.
\end{proof}

\begin{theorem}[Successful-handoff prefix safety]
\label{th:safety}
Assume engine safety, a unique \CutCert, an accepted complete manifest, and no certified abort.
For the admitted locked protocol assume durable \cref{rule:lock} with $q_1>2f$;
the distinct fence-only argument instead assumes a valid \FenceCert{} satisfying
\cref{th:authority-exclusion}. Also assume target initialization and authorization from the
certified root under \cref{def:bootstrap,def:fencecert}. Then correct validators completing the
forward handoff commit one common legacy prefix followed by one target extension from
$r_{\hs-1}$, with no conflicts at any height.
\end{theorem}
\begin{proof}
Legacy safety covers $h<\hd$ and \cref{lem:single} covers $[\hd,\hs)$. Either
\cref{lem:fence-live}(i) or \cref{th:authority-exclusion}, according to the chosen premise,
prevents any competing source commit at $h\ge\hs$, including under partial delivery.
\Cref{th:unique} fixes the parent and root, and \cref{th:unforge} fixes the complete accepted body.
Target initialization starts from $r_{\hs-1}$ only after the local manifest and fence are durable, and
target safety then yields a single canonical extension.
\end{proof}

\begin{theorem}[Conditional completion]
\label{th:liveness}
After GST, suppose every correct validator is eligible to attest and ready on the same fixed
candidate $\beta$, can sign the same valid manifest body, and remains available through handoff.
Assume \cref{rule:lock}, an admitted pair $(q_1,q_F)$ with $q_1>2f$, $q_F\le n-f$ and
$q_F+q_1>n+f$, $\hs<\hr$, bounded local verification and durable writes, and certificate collection
and delivery bounds for every stage, including collector failover and retransmission. Correct
validators durably adopt the certified manifest before installing and sharing its fence, and no
certified abort interrupts the forward handoff. Finally,
suppose the initialized target supplies a finite bound $T_{\mathrm{first}}$ to its first finalization.
Then the forward handoff reaches a first target commit. Let $T_{\mathrm{ready}}$ bound the time from
the last source commit to common eligibility and readiness, including boundary delivery and shadow
execution. If this additional bound exists, the source-to-first-target pause satisfies
\begin{equation}
  \begin{aligned}
  T_{\mathrm{bdy}}\le{}& T_{\mathrm{ready}}+\Delta_B+\Delta_M+\Delta_F\\
  &+T_B(n)+T_M(n)+T_F(n)+T_{\mathrm{init}}+T_{\mathrm{first}},
  \end{aligned}
  \label{eq:tbdy}
\end{equation}
Here $\Delta_{\{B,M,F\}}$ bound collection and delivery of each certificate;
$T_{\{B,M,F\}}(n)$ bound local signing, verification, and durable stage work; and
$T_{\mathrm{init}}$ bounds durable target initialization. The target-specific
$T_{\mathrm{first}}$ is an adapter obligation, not a universal two-message-delay bound for BFT
targets. This theorem assumes common-candidate convergence and successful certificate transport;
it proves neither of them. It also does not bound time before GST, indefinite crashes, incompatible
candidate locks, retry, or terminal abort/seal convergence. If a partition blocks any certificate,
activation stays pending; healing helps only if the same tuple remains viable and its shares are
retained or retransmitted.
\end{theorem}
\begin{proof}
Correct attestations name the immutable $\beta$, so $n-f$ of them form the unique \CutCert{}
(\cref{th:unique}) within the assumed collection/delivery bound. All correct validators accept and
sign the same manifest, which forms the \ManCert{} under its transport bound.
Each correct validator holds the delivered $\pi_{\mathrm B}$ and verifies and durably adopts
$\mathcal M_{\mathrm B}$ within the local stage bounds. These facts discharge the premises of
\cref{lem:fence-live}(ii): the $n-f\ge q_F$ correct validators can install fences and supply enough
shares without Byzantine cooperation. Fence transport delivers their certificate, and the target's
assumed $T_{\mathrm{first}}$ bound yields its first commit. Summing the sequential stage costs
gives \eqref{eq:tbdy}.
\end{proof}

\subsection{Rollback and manifest integrity}
\begin{lemma}[Replay equivalence]
\label{lem:replay}
Under \cref{as:exec,as:metadata}, replaying the ordered transactions of a discarded suffix
$[\hs,h_a]$ onto $\Csrc^{g_1'}$ from $X_{\hs-1}$ reproduces the target attempt's state sequence and
roots. Receipts and events outside $\delta_{\mathrm{exec}}$ are not claimed. If \cref{as:metadata}
fails, equivalence is not guaranteed, and missing or mismatched bound context makes automatic replay
fail closed.
\end{lemma}

\begin{theorem}[Rollback-anchor integrity]
\label{th:rollback}
A correct validator that performs anchor-based restoration at $h_a<\hr$ discards only its provisional
suffix, restores $X_{\hs-1}$, reverts no absolutely final block, and rejects any resume state whose
root differs from $r_{\hs-1}$. A valid $n-f$ \AbortCert{} binding $g_1'>g_1$ separately gates the
committee-scoped transition to $\Csrc^{g_1'}$, and $g_1$ stays fenced. These properties hold without
\cref{as:metadata}.
\end{theorem}

\begin{theorem}[Complete-manifest integrity and agreement]
\label{th:unforge}
Under \cref{as:crypto}, local boundary verification, and one manifest share per correct validator per
epoch: (i)~no adversary controlling at most $f$ keys produces an accepted body whose boundary differs
from the locally finalized state, except with negligible probability; and (ii)~two valid \ManCert{}s
accepted by correct validators authenticate the same body.
\end{theorem}

\begin{proposition}[Replay resistance]
\label{prop:replay}
A manifest valid for $(\mathrm{cid},e)$ verifies for $(\mathrm{cid}',e')\neq(\mathrm{cid},e)$ only with
negligible probability, since both values are inside every signed message.
\end{proposition}

\ref{app:proofs} gives the proofs of \cref{lem:replay,th:rollback,th:unforge}. The
rollback-cycle accounting is conditional on a per-epoch retry cap $\rho$ and a progress ticker. Under
these premises, the time charged to forced rollbacks is at most $\rho(\tau+T_{\mathrm{recover}})$ per
epoch. This bounds only the charged cycles, not a stall that persists after the cap.

\subsection{Shadow detection is coverage-bounded}
By \cref{as:exec}, the shadow root is a deterministic function of its inputs. A root-changing
divergence on an exercising block is detected by a correct validator that performs shadow validation,
not by a per-block random trial. Let $p_{\mathrm{div}}$ be the probability of a latent divergence
under a specified workload distribution, and let $u(\kappa)$ be the conditional probability,
given that divergence, that no block in a length-$\kappa$ window exercises it. Then
\begin{equation}
  \Pr[\text{latent divergence unexercised in window}]=p_{\mathrm{div}}\,u(\kappa),
  \label{eq:coverage}
\end{equation}
where $u$ is non-increasing for nested windows under the same workload process. This is a
coverage identity, not a measured cutover-failure probability: a path that is never exercised has
$u(\kappa)=1$ for every $\kappa$. An independent-trial decay factor $(1-q)^\kappa$ would require
independent path-exercising events, which deterministic replication and shared state do not
guarantee.



\begin{table}[t]
\centering
\caption{Measured single-host Ed25519 share-list cost. ``Sign'' is sequential aggregate CPU time;
size excludes certificate-body and transport framing. Only aggregate timings were retained.}
\label{tab:certtiming}
\small
\renewcommand{\arraystretch}{1.14}
\rowcolors{2}{rowtint}{white}
\begin{tabular}{@{}rrrrr@{}}
\toprule
$n$ & $n-f$ & Sign ($\mu$s) & Verify p50/p95 ($\mu$s) & Size (B) \\
\midrule
\input{tables/certificate_rows.tex}
\end{table}

\begin{table*}[t]
\centering
\caption{Evidence by layer. Theorems assume the stated engine and durable-store contracts;
bounded checks do not establish their runtime conformance. ``Component'' denotes a Rust unit or
integration test; ``campaign'' denotes a reported experiment.}
\label{tab:status}
\small
\renewcommand{\arraystretch}{1.16}
\rowcolors{2}{rowtint}{white}
\begin{tabularx}{\textwidth}{@{}L{0.17\textwidth}L{0.22\textwidth}YL{0.24\textwidth}@{}}
\toprule
\textbf{Element} & \textbf{Proved / model-checked} & \textbf{Implemented} & \textbf{Executed in campaigns} \\
\midrule
Boundary \CutCert & \cref{th:window,th:unique}; TLC, $n\in\{4,7,10\}$, all binary splits & Ed25519 $n-f$ certificate; shared gate predicate & RQ1 (simulator); RQ2 (processes) \\
Cross-engine exclusion & \cref{th:authority-exclusion,th:fence-admission}; TLC partial-delivery, $n\in\{4,5\}$, four falsifying controls & Admission predicate in governance validator and \FenceCert{} verifier; per-validator adopt$\to$install$\to$share$\to$verify order; restart reload & Local fence persistence and restart on the PoA path; committee-wide \FenceCert{} formation over the network is \emph{not} executed \\
Complete manifest & \cref{th:unforge}; TLC, $n=4$ & \ManCert{} verification; durable one-body-per-epoch slot & Manifest shares are not disseminated by the network path \\
Attestation lock & \cref{lem:fence-live} proves quiescence assuming Rule~\ref{rule:lock}; persistence not model-checked & Attestation slot and source prohibition are specified; full runtime conformance not established & No crash-point campaign for the complete Rule~\ref{rule:lock} invariant \\
Abort, seal, rollback & \cref{th:rollback}; TLC rollback model & \AbortCert/\SealCert, terminal lock, local restoration, off-path progress ticker, fail-closed context checks & Context admission (RQ4); no end-to-end certified abort \\
Client finality tier & -- & Tier inside the signed finality payload; effective promotion requires a separate seal proof & Receipts in process campaigns; certified client reconciliation not exercised \\
Raft source & Interface-level (\cref{def:engine}) & Raft adapter with internal fence rejection & Component crossing into HotStuff only \\
\bottomrule
\end{tabularx}
\end{table*}

\begin{figure*}[t]
\centering
\includegraphics[width=\textwidth]{figures/fig_results.pdf}
\caption{Evaluation summary. (a)~Process controls separate unsafe-threshold conflicts from
unauthorized crossings. (b)--(d)~Simulator cutover-evidence latency by seed, committee size, and
one-way delay; these panels report simulator time, not deployment latency. Whiskers in (c)--(d) show
retained seed ranges, not confidence intervals.}
\label{fig:results}
\end{figure*}

\subsection{Complexity and certificate cost}
Shadow validation adds no wire messages. With an available collector and bounded delivery to all
validators, each of the three activation-certificate stages takes $O(n)$ share submissions and
$O(n)$ certificate deliveries over $O(1)$ post-GST communication rounds. These design-level
counts exclude unbounded retries and assume the collector/failover and local-work bounds of
\cref{th:liveness}; they are not measurements of the networked path.
All-to-all share dissemination instead costs $O(n^2)$ messages per stage. Distributing an
unaggregated $O(n)$-signature certificate to $n$ validators costs $O(n^2)$ signature words.
During dual-run each validator performs one off-path re-execution per block, in $O(|B_h|)$
time, and stores $O(\kappa)$ verdicts, a constant number of $O(n)$ certificates, and at most $\hr-\hs$
provisional blocks. \Cref{tab:certtiming} reports measured Ed25519 costs. Each share is $72$\,B
(a $64$-byte signature and an $8$-byte signer index), so the share list occupies $72(n-f)$ bytes.
This excludes the signed body, collection encoding, and transport framing; it is not a wire-size
measurement. At $n=200$ the $134$ shares verify in $3.9$\,ms at the median and their payload is
under $10$\,KiB.

